%\authorRemark{Notes: ensure things are not in first person}

\chapter{Introduction}
\label{chap:introduction}

While there have been impressive advances in artificial intelligence in recent years, modern AI systems remain fundamentally flawed, lacking transparency and true reasoning capabilities \cite{jiang2024peektokenbiaslarge}. Despite breakthroughs in natural language processing \cite{nlp_deep_learning_survey}, computer vision \cite{cv_deep_learning_survey}, and reinforcement learning \cite{rl_deep_learning_survey}, deep learning models continue to operate as opaque black boxes, offering little transparency into how or why decisions are made. Furthermore, while recent work on chain-of-thought has improved model performance \cite{wei2023chainofthoughtpromptingelicitsreasoning}, it gives the illusion that large language models are capable of actual ``reasoning'' through genuine logical inference and generalization rather than relying solely on statistical pattern matching \cite{bastounis2024consistentreasoningparadoxintelligence, li202512surveyreasoning, cherian2023deepneuralnetworkssmarter}. As a result, AI remains untrustworthy in high-stakes environments that deal with human lives, as it is prone to hallucinations \cite{maleki2024aihallucinationsmisnomerworth}, biases \cite{simhi2025trustmeimwrong, echterhoff2024cognitivebiasdecisionmakingllms}, and inexplicable errors \cite{rogha2023explaindecidehumancentricreview}. Without an explicit mechanism for structured reasoning and transparency, deep learning will continue to encounter significant challenges in environments that demand reliability, interpretability, and real-time decision making.

This current lack of transparency in AI systems gives way to fields such as Explainable Artificial Intelligence (XAI), which has become increasingly important as AI systems are deployed in high-stakes and decision-critical domains. Black-box models, especially deep neural networks with millions of parameters, often lack transparency, making it difficult for users to trust their decisions. The AI community recognizes that explainability is a crucial feature for the practical deployment of AI models, as it helps overcome the opaqueness and complexity presented by subsymbolic learning techniques and ensure AI models deployed in critical sectors can be trusted.
In response, XAI research has produced methods to interpret model predictions, ranging from visualization of learned features to extraction of human-understandable rules \cite{interpreting-nns}. These methods aim to build user trust by shedding light on how inputs affect outputs, thus moving towards Responsible AI \cite{baker2023explainableairesponsibleai}. However, post hoc (techniques applied after a model is trained) XAI methods remain fundamentally flawed, as they attempt to interpret models that are inherently uninterpretable \cite{inscrutable_xai_algorithms}. 

% Studies have shown that such techniques often provide misleading or incomplete explanations \cite{inscrutable_xai_algorithms}, reinforcing the need for intrinsic interpretability.

A promising approach to address these limitations is Neurosymbolic AI (NSAI) \cite{sarker2021neurosymbolicartificialintelligencecurrent}. NSAI integrates subsymbolic (neural) and symbolic paradigms, combining data-driven machine learning with structured symbolic reasoning. Neural networks excel at recognizing patterns but remain opaque, while symbolic systems provide logical clarity but struggle with adaptability. By merging these two approaches, NSAI aims to create AI systems that are both inherently interpretable and scalable \cite{zhang2024neurosymbolicaiexplainabilitychallenges, sheth2023neurosymbolicaiwhy}. 

This thesis focuses on a type of NSAI architecture called Differentiable Logic Gate Networks \cite{difflogic} and develops upon it through the following novel contributions:

\begin{enumerate}[noitemsep,topsep=0pt]

     \item \textbf{FPGA-Based Deployment of DiffLogic:} Implemented DiffLogic models on Field Programmable Gate Arrays (FPGAs), demonstrating real-time inference capabilities that significantly increase speeds.
    \item \textbf{MiniNet Compression Algorithm:} Developed a method for reducing the size of trained Differentiable Logic Gate Networks while preserving their accuracy and improving efficiency.
    \item \textbf{Visualization Interface for Interpretability:} Designed a tool that visualizes the logic structures learned by DiffLogic and a separate tool that generates saliency maps with different XAI algorithms for comparison, making model decisions more accessible and understandable for human users.
\end{enumerate}

To structure this exploration, Chapter \ref{chap:background} goes over limitations of current AI systems, arguments against purely connectionist models, the promise of neurosymbolic approaches, and a detailed explanation of DiffLogic Networks. Chapter \ref{chap:methods} outlines the methods, including experiments run using the base DiffLogic network, hardware deployment via FPGAs, compression into MiniNets, and the development of visualization tools. Chapter \ref{chap:results} presents experimental results and comparisons across all approaches. Chapter \ref{chap:conclusion} concludes the thesis by summarizing key findings, discussing limitations, and proposing future research directions, particularly in expanding interpretability and efficiency.
